%% =====================================================================
%%  B1-T-19-03 -- what B1 contributes to "The Unanswerable Question"
%%  -------------------------------------------------------------------
%%  LAYER A: APPEND-ONLY. Written once at the entry's write, then left
%%  alone. If the source has more to say later, add a bullet -- do not
%%  rewrite. Material found ONLY here goes in the `unique` slot with
%%  @unique set to yes, which makes it undeletable (rule R10).
%% =====================================================================
%% @kind: source
%% @id: B1-T-19-03
%% @book: B1
%% @topic: T-19-03
%% @locator: Tactic 12 (ch. XIII), pp. 153--160 (social-pressure section)
%% @status: distilled
%% @unique: no
%% @integrated: yes
%% @updated: 2026-09-19

\dossier{B1}{T-19-03}{\bookshort{B1}}{Tactic 12 (ch. XIII), pp. 153--160 (social-pressure section)}

\begin{distilled}
  \item The strolling minstrel asks what song you like: the question has no grammatical room for no, and the watching audience converts it into a dollar.
  \item B1 classes social pressure as a working type of pressure: arrange the scene so declining costs the target's public self-image more than compliance costs their wallet.
  \item The device's cousins in the same chapter: the assumed yes and the silence after the ask --- all work by pricing the target's exit above their consent.
  \item The chapter's own tell: the author resented the dollar and enjoyed the song none --- compliance under audience pressure leaves resentment, not agreement.
\end{distilled}
